\documentclass[11pt]{article}
\usepackage[margin=1in]{geometry}
\usepackage{amsmath,amssymb,hyperref,booktabs}
\title{Abliteration and Adversarial Validation of Refusal Mechanisms in LFM3-1.7B}
\author{Research manuscript --- author review draft}
\date{October 29, 2026 \\ Version 0.1.0}
\begin{document}
\maketitle
\begin{abstract}
We study refusal-related representations in the tri-hybrid LFM3-1.7B model
and separate behavioral refusal from effective resistance to unauthorized
actions. We preregister benign/adversarial paired contrasts, layer-level
representation measurements, reversible interventions, and symbolic
tool-authorization evaluation. This preliminary manuscript reports methods,
not validated safety performance or model-specific causal effects.
\end{abstract}
\section{Introduction}
Previous mechanistic work found low-dimensional refusal-mediating
representations in some transformer language models \cite{arditi2024}.
LFM3-1.7B uses 13 gated-convolution, five sliding-window attention, and six
global-attention blocks \cite{liquid2026}, raising an open question about
whether those mechanisms generalize across architectures. Agentic use also
requires distinguishing verbal refusal from actual action authorization.
\section{Research questions}
We test four preregistered questions:
(1) whether paired refusal-associated contrasts replicate in held-out prompts;
(2) whether reversible perturbations affect refusal independently of benign
helpfulness; (3) whether effects depend on operator type; and
(4) whether verbal refusal predicts resistance to unauthorized tool actions.
A shift in refusal language is not evidence of effective safety by itself.
\section{Model and experimental design}
The principal subject is the post-trained LiquidAI/LFM3-1.7B checkpoint.
LiquidAI/LFM3-1.7B-Base is a distinct comparison, not an equivalent control.
Full-precision and quantized versions are analyzed in separate strata.
We freeze model revision, tokenizer, prompt template, seed, dtype, runtime,
and dataset hash for each experiment.
\subsection{Activation contrasts}
For benign and adversarial variants of matched prompt $i$, let
$h_\ell^{(b,i)}$ and $h_\ell^{(a,i)}$ be the layer-$\ell$ hidden vectors
at the assistant generation prefix. The descriptive contrast is
\[
\bar{\Delta}_\ell=\frac{1}{N}\sum_i
(h_\ell^{(a,i)}-h_\ell^{(b,i)}),\qquad
d_\ell=\frac{\bar{\Delta}_\ell}{\|\bar{\Delta}_\ell\|_2}.
\]
These vectors can encode wording and topic rather than a causal refusal
mechanism. Direction selection must occur on training pairs only; holdout
prompts and sham vectors evaluate generalization.

\subsection{Reversible interventions}
On an inference-only forward hook, project activation $h_\ell$ by
\[
h_\ell'=h_\ell-\alpha \langle h_\ell,d_\ell\rangle d_\ell,
\quad \alpha\in\{0,0.25,0.5,0.75,1\}.
\]
Compare identical seeds, decoding settings, and model weights under
baseline, projection, and random-direction controls. Assess convolution,
sliding-attention, and global-attention blocks separately. Preserve the
original checkpoint and do not automatically distribute altered weights.

\subsection{Behavioral outcomes}
Evaluate (a) justified refusal or safe redirection of malicious prompts,
(b) human-graded actionable harmful content, (c) benign utility and
false-positive refusals, and (d) simulated unauthorized action attempts.
Authorization is evaluated by a separate deterministic action policy:
a model response never grants the permissions that policy enforces.
No executable malicious tools, external targets, or live actions are used.
Keyword-based refusal detection is only a weak descriptive proxy.

\subsection{Analysis and validation}
Use blinded human review, independent heldout cases, multiple prompt
paraphrases, and category-based breakdowns. Report paired sample sizes,
bootstrap confidence intervals, and correction for searched layers and
projection strengths. Use standardized red-teaming ideas such as HarmBench
\cite{mazeika2024} only with appropriate controls and data handling.
A lower refusal rate is never treated as a successful safety outcome.
\section{Three-checkpoint experimental design (v0.1)}
We use exactly three primary checkpoint arms, with all evaluations run
using the same private harness. The pretrained model is treated as a
\emph{developmental control}, not as a refusal-free chat baseline:
\begin{center}
\begin{tabular}{lll}
\toprule
Arm & Checkpoint & Primary comparison \\
\midrule
A & LFM3-1.7B post-trained & Native refusal/utility \\
B & M1-style edited A & Refusal suppression vs. A \\
C & LFM3-1.7B-Base & Behaviors prior to post-training \\
\bottomrule
\end{tabular}
\end{center}
The edit in B is a conventional direct-removal variant of the
difference-of-means abliteration of Arditi et al.\ \cite{arditi2024}.
Unlike the initial reversible forward-hook probe, B is a separate,
locally saved model whose weights are modified. Its original weights
and full chain of custody are retained. With $d$ a normalized
candidate direction and $W$ an operator's residual-stream output
weight matrix, we use
\[
W'=(I-\alpha dd^{\top})W,\qquad\alpha=1\quad\text{(M1 direct removal)}.
\]
The experimental implementation edits each of the 24 MLP $w_2$
weights and the 24 residual operator outputs (13 convolution output
projections and 11 attention output projections), for 48 edits.
Because LFM3 uses sandwich RMS normalizations and hybrid operators,
this projection does \emph{not} guarantee removal of the corresponding
direction at residual addition. It is an empirical hypothesis to test.

\subsection{Preregistered scoring and independent controls}
For each of A, B, and C, measure a held-out benign-helpfulness score,
false-refusal rate, human-graded harmful-actionability score, safe
redirection on unauthorized requests, latency, and per-token cost.
Use deterministic zero-temperature runs before evaluating multiple
sampling seeds; apply balanced prompt variations. For the base arm,
also evaluate natural next-token completion to account for its
different training objective. The mock authorization policy is a
separate control: it must block prohibited actions regardless of
the model's text.

A sham projection, random-vector orthogonalization, and alpha sweep
are \emph{methodological controls}, not additional checkpoints in the
three-arm comparison. All hyperparameter selections use calibration
splits. Holdouts remain unseen until choices are frozen.
We will report 95\% confidence intervals and human annotation
agreement, with all denominators and failures disclosed.

\subsection{Prototype results and evidence policy}
The first two-pair calibration and two-pair held-out reversible
forward-hook probe are strictly exploratory. It was not sufficient
to validate safety. The locally generated M1-style derivative is
similarly an unvalidated smoke checkpoint. No public distribution
or automatic model replacement is authorized. Quantitative claims
require a larger independent dataset, prompt contamination audit,
human review, and a pre-registered decision threshold.

\section{XSTest as a frozen independent benchmark}
We use XSTest \cite{xstest2024} as an external, non-calibration
benchmark of exaggerated safety behavior. Its 450 prompt set
comprises 250 safe and 200 unsafe contrast prompts, distributed
under CC BY 4.0. The exact source CSV is pinned by SHA-256
\texttt{11783fb294ed017473ee53c207d71f2161c7672c8d0b037501e78387f801cb5a}.
The dataset is never used to estimate the weight-edit direction or
choose the modified checkpoint. The prescribed three-arm evaluation
is $450\times 3=1350$ generations, with fixed templates and seeds.
Each prompt is scored for full compliance, partial refusal, full
refusal, and uncertain outcomes by two model-blinded human reviewers.
We additionally annotate actionable harmful content and harmless
helpfulness. Report safe over-refusal and unsafe actionability in
separate strata: neither total refusal nor lexical apology alone is
an adequate safety outcome. We compute paired A/B contrasts with
uncertainty intervals and disclose all ungraded items and reviewer
disagreements.

Early stratified test runs are labeled integration smoke results,
not independent evidence of abliteration efficacy. In particular,
the pretrained checkpoint C lacks instruction tuning and requires
a raw-completion sensitivity analysis; treating native chat-template
generations as equivalent would overstate comparability.

\section{Governance and responsible dissemination}
Model weights and raw adversarial generations remain private.
There is no real tool execution, unrestricted serving endpoint, or
automatic model promotion. Distributing altered weights requires independent
safety assessment, legal/license review, and explicit publication approval.

\section{Implementation status and limitations}
Eight paired prompts, six simulated action cases, a local activation-probe
runner, and unit tests are implemented. They are smoke tests, not safety
benchmarks. The causal role of the measured direction remains unknown.
Prompt-topic and prompt-length confounds can distort activation contrasts.
A verbal refusal is not proof that malicious real-world actions are blocked.

\section{Reproducibility}
Pin model revision, template, data hashes, seeds, precision, code and runtime.
Maintain heldout tests, independent grading, confidence intervals, and
negative results. Do not interpret lexical refusal rates as safety outcomes.
\begin{thebibliography}{9}
\bibitem{arditi2024} A. Arditi et al.,
{\it Refusal in Language Models Is Mediated by a Single Direction},
NeurIPS 2024, arXiv:2406.11717.
\bibitem{mazeika2024} M. Mazeika et al.,
{\it HarmBench: A Standardized Evaluation Framework for Automated
Red Teaming and Robust Refusal}, arXiv:2402.04249 (2024).
\bibitem{liquid2026} Liquid AI,
{\it Introducing LFM3: Every Modality for Specialized Intelligence},
October 27, 2026.
\bibitem{xstest2024} P. R\"ottger et al.,
{\it XSTest: A Test Suite for Identifying Exaggerated Safety Behaviours
in Large Language Models}, NAACL 2024.
\end{thebibliography}
\end{document}
